\chapter{Conclusion}\label{ch10:conclusion}

This study evaluates the performance of transformer-based architectures and traditional regression models for time-series forecasting across multiple datasets and forecasting horizons. The results indicate that transformer-based models consistently outperform regression baselines, particularly for long-term forecasting.

Among the models evaluated, \textbf{SAMFormer} demonstrates the most stable and reliable performance across all datasets and forecasting horizons. It maintains high accuracy and exhibits minimal performance degradation as the forecasting horizon increases. \textbf{TimeGPT} achieves the best results in short-term forecasting, effectively capturing immediate temporal dependencies, but its accuracy declines for longer horizons. \textbf{Time-MoE} performs competitively in certain domains, particularly in healthcare forecasting, where irregular temporal patterns present challenges for other models. However, its effectiveness varies across datasets. In contrast, \textbf{TimesFM} underperforms across multiple domains, showing limited improvements even after fine-tuning.

The results also highlight the impact of fine-tuning on forecasting accuracy. Fine-tuning leads to significant performance gains for \textbf{TimeGPT} and \textbf{Time-MoE}, particularly for longer horizons, with reductions in forecasting errors of up to 15\%. \textbf{TimesFM}, however, does not show notable improvements from fine-tuning, maintaining lower accuracy compared to other transformer-based models. Across all foundation models, fine-tuning results in lower forecasting errors compared to zeroshot performance, reinforcing the importance of adapting models to specific datasets.

Traditional regression models, including \textbf{Linear Regression} and ensemble-based methods, perform well for short-term forecasting but struggle with longer horizons. Their limited capacity to model complex temporal dependencies leads to significant performance degradation over extended prediction windows. While these models remain viable for applications requiring low computational complexity and short-term predictions, they are less effective for tasks demanding robust long-term forecasting capabilities.

Overall, the findings confirm that transformer-based architectures offer significant advantages over traditional methods for time-series forecasting, particularly in long-horizon settings. However, no single model achieves superior performance across all datasets, underscoring the need for dataset-specific fine-tuning. Future research should explore hybrid approaches that combine the efficiency of regression-based methods with the predictive power of transformers. Further optimization of \textbf{TimeGPT} could improve its long-term forecasting performance, while strategies to enhance computational efficiency may increase the accessibility of transformer-based models for real-world applications. Continued work in refining fine-tuning methodologies and expanding training datasets could further improve forecasting accuracy, ensuring that these models can be effectively applied across a broad range of time-series prediction tasks.
